\documentclass[10pt,a4paper]{amsart}
\usepackage{amsmath,amssymb,amsthm,graphicx}
\usepackage{verbatim}
\newenvironment{reference}{\comment}{\endcomment}
%\newenvironment{reference}{}{}
\newenvironment{slogan}{\comment}{\endcomment}
\newenvironment{history}{\comment}{\endcomment}

% For commutative diagrams we use Xy-pic
\usepackage[all]{xy}

% We use 2cell for 2-commutative diagrams.
\xyoption{2cell}
% Theorem environments.
%
\theoremstyle{plain}
\newtheorem{theorem}[subsection]{Theorem}
\newtheorem{proposition}[subsection]{Proposition}
\newtheorem{lemma}[subsection]{Lemma}

\theoremstyle{definition}
\newtheorem{definition}[subsection]{Definition}
\newtheorem{example}[subsection]{Example}
\newtheorem{exercise}[subsection]{Exercise}
\newtheorem{situation}[subsection]{Situation}

\theoremstyle{remark}
\newtheorem{remark}[subsection]{Remark}
\newtheorem{remarks}[subsection]{Remarks}

\numberwithin{equation}{subsection}

% Macros
%
\def\lim{\mathop{\mathrm{lim}}\nolimits}
\def\colim{\mathop{\mathrm{colim}}\nolimits}
\def\Spec{\mathop{\mathrm{Spec}}}
\def\Hom{\mathop{\mathrm{Hom}}\nolimits}
\def\Ext{\mathop{\mathrm{Ext}}\nolimits}
\def\SheafHom{\mathop{\mathcal{H}\!\mathit{om}}\nolimits}
\def\SheafExt{\mathop{\mathcal{E}\!\mathit{xt}}\nolimits}
\def\Sch{\mathit{Sch}}
\def\Mor{\mathop{\mathrm{Mor}}\nolimits}
\def\Ob{\mathop{\mathrm{Ob}}\nolimits}
\def\Sh{\mathop{\mathit{Sh}}\nolimits}
\def\NL{\mathop{N\!L}\nolimits}
\def\CH{\mathop{\mathrm{CH}}\nolimits}
\def\proetale{{pro\text{-}\acute{e}tale}}
\def\etale{{\acute{e}tale}}
\def\QCoh{\mathit{QCoh}}
\def\Ker{\mathop{\mathrm{Ker}}}
\def\Im{\mathop{\mathrm{Im}}}
\def\Coker{\mathop{\mathrm{Coker}}}
\def\Coim{\mathop{\mathrm{Coim}}}

% Boxtimes
%
\DeclareMathSymbol{\boxtimes}{\mathbin}{AMSa}{"02}

%
% Macros for moduli stacks/spaces
%
\def\QCohstack{\mathcal{QC}\!\mathit{oh}}
\def\Cohstack{\mathcal{C}\!\mathit{oh}}
\def\Spacesstack{\mathcal{S}\!\mathit{paces}}
\def\Quotfunctor{\mathrm{Quot}}
\def\Hilbfunctor{\mathrm{Hilb}}
\def\Curvesstack{\mathcal{C}\!\mathit{urves}}
\def\Polarizedstack{\mathcal{P}\!\mathit{olarized}}
\def\Complexesstack{\mathcal{C}\!\mathit{omplexes}}
% \Pic is the operator that assigns to X its picard group, usage \Pic(X)
% \Picardstack_{X/B} denotes the Picard stack of X over B
% \Picardfunctor_{X/B} denotes the Picard functor of X over B
\def\Pic{\mathop{\mathrm{Pic}}\nolimits}
\def\Picardstack{\mathcal{P}\!\mathit{ic}}
\def\Picardfunctor{\mathrm{Pic}}
\def\Deformationcategory{\mathcal{D}\!\mathit{ef}}
\usepackage{cleveref}
\crefname{originallemma}{original lemma}{original lemmas}
\Crefname{originallemma}{Original lemma}{Original lemmas}
\makeatletter
\newlabel{formal-spaces-lemma-characterize-quasi-compact-morphism}{{0AJB}{source}{}{sourcepointer-1}{}}
\newlabel{formal-spaces-lemma-characterize-quasi-compact-morphism@cref}{{[originallemma][0][]0AJB}{[1][0][]source}{}{}{}}
\makeatother
\crefname{originalreference}{original reference}{original references}
\Crefname{originalreference}{Original reference}{Original references}
\makeatletter
\newlabel{lemma-functoriality-completion-functor}{{0GDQ}{source}{}{sourcepointer-2}{}}
\newlabel{lemma-functoriality-completion-functor@cref}{{[originalreference][0][]0GDQ}{[1][0][]source}{}{}{}}
\makeatother
\crefname{originalreference}{original reference}{original references}
\Crefname{originalreference}{Original reference}{Original references}
\makeatletter
\newlabel{lemma-quasi-compact-gives-quasi-compact}{{0GDE}{source}{}{sourcepointer-3}{}}
\newlabel{lemma-quasi-compact-gives-quasi-compact@cref}{{[originalreference][0][]0GDE}{[1][0][]source}{}{}{}}
\makeatother
\crefname{originallemma}{original lemma}{original lemmas}
\Crefname{originallemma}{Original lemma}{Original lemmas}
\makeatletter
\newlabel{spaces-morphisms-lemma-quasi-compact-local}{{03KG}{source}{}{sourcepointer-4}{}}
\newlabel{spaces-morphisms-lemma-quasi-compact-local@cref}{{[originallemma][0][]03KG}{[1][0][]source}{}{}{}}
\makeatother
\crefname{originalreference}{original reference}{original references}
\Crefname{originalreference}{Original reference}{Original references}
\makeatletter
\newlabel{theorem-dilatations-general}{{0ARB}{source}{}{sourcepointer-5}{}}
\newlabel{theorem-dilatations-general@cref}{{[originalreference][0][]0ARB}{[1][0][]source}{}{}{}}
\makeatother
\title{Quasi-compactness transfers between a dilatation and its formal model}
\author{The Stacks Project Authors}
\date{stacks-project:restricted.tex:local-snapshot-20261012; GNU FDL 1.2 or later}
\begin{document}
\maketitle
\noindent\textit{Editorial note.} All delivered mathematical text is the author's own, copied from the pinned Stacks Project chapter file restricted.tex; no statement or proof step was rewritten, reordered or removed.\par\medskip
\noindent\textit{Editorial note.} The counted proof is the author's complete argument: the proof reduces etale locally to the affine case through the functoriality of the completion functor, extracts a finite subcover of the formal part from quasi-compactness of W, and covers X' by the open part plus the finitely many pieces mapping onto W\_red\par\medskip
\noindent\textit{Editorial note.} The label lemma-output-quasi-compact names the counted statement, and the proof environment immediately following it in the source is delivered verbatim as its proof.\par\medskip
\noindent\textit{Editorial note.} The counted statement is carried by the author's own lemma environment in the chapter source.\par\medskip
\noindent\textit{Allowed standard prerequisites (editorial).}\begin{itemize}
\item the dilatations theorem (theorem-dilatations-general of the same chapter) and its correspondence between f : X' -> X and g : W -> X\_\{/T\}
\item quasi-compactness of f implies quasi-compactness of g (lemma-quasi-compact-gives-quasi-compact of the same chapter)
\item quasi-compactness of morphisms of algebraic spaces checked etale locally on the base (Morphisms of Spaces, Lemma spaces-morphisms-lemma-quasi-compact-local)
\item quasi-compactness of morphisms of formal algebraic spaces characterized etale locally (Formal Spaces, Lemma formal-spaces-lemma-characterize-quasi-compact-morphism) and the functoriality of the completion functor (lemma-functoriality-completion-functor of the same chapter)
\item surjectivity of etale covers and the description |X'| = |U| amalgramated with |W\_red| for the dilatation
\end{itemize}\medskip
\ifdefined\hypertarget\hypertarget{sourcepointer-1}{}\fi
\noindent\textit{Source reference.} Original lemma 0AJB: sources/stacks-chow/tags.txt, source lines 13294--13294. Prerequisite pointer invoked at proof line(s) 7132; it names a standard fact of a sibling chapter that the counted argument uses, and every step of the counted proof built on it is retained verbatim in the delivered body.\par\medskip
\ifdefined\hypertarget\hypertarget{sourcepointer-2}{}\fi
\noindent\textit{Source reference.} Original reference 0GDQ: sources/stacks-restricted/pinned-original-source.tex, source lines 6790--6790. Prerequisite pointer invoked at proof line(s) 7135; it names a standard fact of the same chapter that the counted argument uses, and every step of the counted proof built on it is retained verbatim in the delivered body.\par\medskip
\ifdefined\hypertarget\hypertarget{sourcepointer-3}{}\fi
\noindent\textit{Source reference.} Original reference 0GDE: sources/stacks-restricted/pinned-original-source.tex, source lines 6067--6067. Prerequisite pointer invoked at proof line(s) 7125; it names a standard fact of the same chapter that the counted argument uses, and every step of the counted proof built on it is retained verbatim in the delivered body.\par\medskip
\ifdefined\hypertarget\hypertarget{sourcepointer-4}{}\fi
\noindent\textit{Source reference.} Original lemma 03KG: sources/stacks-chow/tags.txt, source lines 4676--4676. Prerequisite pointer invoked at proof line(s) 7130; it names a standard fact of a sibling chapter that the counted argument uses, and every step of the counted proof built on it is retained verbatim in the delivered body.\par\medskip
\ifdefined\hypertarget\hypertarget{sourcepointer-5}{}\fi
\noindent\textit{Source reference.} Original reference 0ARB: sources/stacks-restricted/pinned-original-source.tex, source lines 7015--7015. Prerequisite pointer invoked at proof line(s) 7118; it names a standard fact of the same chapter that the counted argument uses, and every step of the counted proof built on it is retained verbatim in the delivered body.\par\medskip
\begin{lemma}
\label{lemma-output-quasi-compact}
With assumptions and notation as in Theorem \ref{theorem-dilatations-general}
let $f : X' \to X$ correspond to $g : W \to X_{/T}$.
Then $f$ is quasi-compact if and only if $g$ is quasi-compact.
\end{lemma}

\begin{proof}
If $f$ is quasi-compact, then $g$ is quasi-compact by
Lemma \ref{lemma-quasi-compact-gives-quasi-compact}.
Conversely, assume $g$ is quasi-compact.
Choose an \'etale covering $\{X_i \to X\}$ with $X_i$ affine.
It suffices to prove that the base change $X' \times_X X_i \to X_i$
is quasi-compact, see
Morphisms of Spaces, Lemma \ref{spaces-morphisms-lemma-quasi-compact-local}.
By Formal Spaces, Lemma
\ref{formal-spaces-lemma-characterize-quasi-compact-morphism}
the base changes $W_i \times_{X_{/T}} (X_i)_{/T}  \to (X_i)_{/T}$
are quasi-compact.
By Lemma \ref{lemma-functoriality-completion-functor}
we reduce to the case described in the next paragraph.

\medskip\noindent
Assume $X$ is affine and $g : W \to X_{/T}$ quasi-compact.
We have to show that $X'$ is quasi-compact.
Let $V \to X'$ be a surjective \'etale morphism
where $V = \coprod_{j \in J} V_j$ is a disjoint union of affines. Then
$V_{/T} \to X'_{/T} = W$ is a surjective \'etale morphism.
Since $W$ is quasi-compact, then we can find a finite subset
$J' \subset J$ such that $\coprod_{j \in J'} (V_j)_{/T} \to W$ is surjective.
Then it follows that
$$
U \amalg \coprod\nolimits_{j \in J'} V_j \longrightarrow X'
$$
is surjective (and hence $X'$ is quasi-compact).
Namely, we have $|X'| = |U| \amalg |W_{red}|$ as $X'_{/T} = W$.
\end{proof}

\end{document}
